\documentclass[11pt,a4paper]{ctexart}
\usepackage[a4paper,margin=1.60cm,top=1.8cm,bottom=1.8cm]{geometry}
\usepackage{../../../00_system/templates/ipara-reading}

% ==============================================================================
% 《The Magic of Summer Reading》单篇精读训练手札
%
% 原文与出处：article.md
% 本文件：只保存这一篇文章的理解、声音与输出训练
% 长期表达定义：02_expressions/
% 长期能力状态：03_capabilities/
% ==============================================================================

\begin{document}

\articleheader
  {The Magic of Summer Reading}
  {How America's Competitive Summer Reading Culture Evolved into a Digital Phenomenon}
  {Spotlight · The View from America}
  {American Culture · Summer Reading · Digital Trends}
  {B2--C1（训练标签）}
  {2026-08-20}

% ------------------------------------------------------------------------------
% 一、双栏精读
% ------------------------------------------------------------------------------
\readingsection{一、双栏精读}

\begin{paracol}{2}

\artpara{P1}{
  At this point in the summer, my favorite thing to do is to read on the beach.
  \kw{There's nothing like a good book to help me escape}{1} when the weather is too hot and humid to enjoy doing anything else.
}

\switchcolumn

\noteitem{1}{There's nothing like... to... · EXP-006}
  {当前语境：没有什么比一本好书更能让作者从炎热潮湿的天气中暂时抽离。}
  {Nothing beats a good book when you just want to switch off.}
  {完整长期节点见表达层；这里保留它在本段中的具体作用。}

\switchcolumn*

\artpara{P2}{
  Beach reading may not be \obs{unique to Americans}{1}, but the American approach to summer reading is different from other countries. It is a \kw{highly organized effort}{2}, promoted nationally by libraries and the media.
}

\switchcolumn

\noteobs{1}{be unique to...}
  {这里不是说美国人独有海滩阅读，而是先否定“行为本身独有”，再转向“美国的组织方式不同”。}

\noteitem{2}{a highly organized effort · 本篇候选}
  {高度组织化、由机构系统推动的活动。}
  {a really well-organized campaign}
  {当前先留在本篇，不因为表达不错就自动建立长期节点。}

\switchcolumn*

\artpara{P3}{
  In Germany, school summer vacation is just six weeks. But in the United States, it can last up to 12 weeks. As a kid, those were the best 12 weeks of the year. But many educators would disagree with my younger self -- they're concerned that children forget too much when they're away from school for such a long period, an issue known as the \kw{``summer slide.''}{3}
}

\switchcolumn

\noteitem{3}{summer slide \ipa{/ˈsʌmə slaɪd/} · 主题术语}
  {长暑假期间出现的知识遗忘或学业退步。这里首先是理解文章背景的主题词。}
  {kids falling behind over the summer}
  {专业术语是否进入长期表达层，取决于未来真实使用需求。}

\switchcolumn*

\artpara{P4}{
  Librarians try to \kw{counter the summer slide}{4} by promoting reading contests. Kids spend the summer months \kw{logging books}{5} that they read and for which they can win prizes, such as coupons for pizza and ice cream. Public figures -- including former President Barack Obama and actress Reese Witherspoon -- also release summer book lists.
}

\switchcolumn

\noteitem{4}{counter the summer slide · EXP-002}
  {通过阅读活动主动抵消暑期学业退步。}
  {stop kids from falling behind over summer break}
  {完整表达节点见表达层。}

\noteitem{5}{log books \ipa{/lɒɡ bʊks/} · 本篇候选}
  {这里的 log 是“记录、登记”，不是名词“原木”。}
  {track what you read}
  {先解决熟词义与搭配；是否长期入库后续再判断。}

\switchcolumn*

\artpara{P5}{
  Some of my dearest childhood memories involve summer reading. After spending the morning at the pool, my brothers and I would often head to the library, where we \kw{curled up in cozy nooks}{6} and read for hours. The summer of 2007 -- when I was nine years old -- was particularly fun. I eagerly awaited the release of the seventh Harry Potter book and \kw{bailed on family plans}{7} for two days in order to read it.
}

\switchcolumn

\noteitem{6}{curl up in a cozy nook · EXP-003}
  {身体蜷坐在舒适的小角落里，带有明显的惬意场景感。}
  {get comfy in a quiet corner with a book}
  {完整长期节点见表达层。}

\noteitem{7}{bail on plans · EXP-001}
  {临时推掉或放弃原本的安排；本段带有轻松口语色彩。}
  {back out of the plans / skip the plans}
  {完整长期节点见表达层。}

\switchcolumn*

\artpara{P6}{
  I'm still an enthusiastic reader today. Like many other American bookworms, I've tried to \kw{keep the magic of summer reading challenges alive}{8} in other ways. My friends and I use the website Goodreads to track the books we've read, share recommendations, and set new reading goals. According to YouGov, Goodreads is the 13th-most-popular social media website among American adults -- ranking higher than TikTok, X, and Snapchat.
}

\switchcolumn

\noteitem{8}{keep the magic of ... alive · EXP-005}
  {让某种原有的乐趣、氛围或传统继续存在。}
  {keep that feeling / tradition going}
  {本段重点是“童年的暑期阅读习惯如何延续到成年”。}

\switchcolumn*

\artpara{P7}{
  I've also recently started using an e-reader. The one I use \kw{lives up to the hype}{9}: It is lightweight, allows for easy annotating, and has backlit functions that are \obs{ideal for sunny days and dark bedrooms alike}{2}. It is also easy to check out e-books virtually from the library.
}

\switchcolumn

\noteitem{9}{live up to the hype · EXP-004}
  {实际体验确实配得上此前的宣传和期待。}
  {it's actually as good as people say}
  {完整长期节点见表达层。}

\noteobs{2}{sunny days and dark bedrooms alike}
  {alike 放在并列对象之后，表示“两种看似相反的场景都一样适用”。}

\switchcolumn*

\artpara{P8}{
  \kw{Amid the boom of BookTok}{10} and Goodreads, \textit{The Wall Street Journal} reported in 2024 that Kindles had become a \kw{``must-have accessory''}{11} among young American women. Sales of the device have gone up in the years since the COVID-19 pandemic. As one \textit{New York Times} author wrote: ``I love reading on it so much, I worry that I may never pick up another physical book again.''
}

\switchcolumn

\noteitem{10}{amid the boom of... · EXP-007}
  {把后文现象放进“某事正在快速兴起”的大背景中。}
  {with ... becoming really popular}
  {完整长期节点见表达层。}

\noteitem{11}{must-have accessory · 本篇候选}
  {从单纯的阅读设备转向“年轻群体的潮流必备品”。}
  {something everyone suddenly wants}
  {当前主要服务本文趋势理解，不急于长期入库。}

\end{paracol}

% ------------------------------------------------------------------------------
% 二、文章结构与主旨
% ------------------------------------------------------------------------------
\readingsection{二、文章结构与主旨复原}

\begin{coretakeaway}[训练后主旨]
  \textbf{American summer reading has evolved from a library-led response to the ``summer slide'' into a broader social reading culture supported by digital platforms and e-readers.}
\end{coretakeaway}

\begin{itemize}
  \item \textbf{个人切入}：作者从自己喜欢暑期阅读的体验开场。
  \item \textbf{制度背景}：美国暑假较长，教育者担心出现 summer slide，图书馆因此组织阅读挑战。
  \item \textbf{个人记忆}：作者用童年阅读、Harry Potter 与家庭经历让制度背景具体化。
  \item \textbf{数字化延续}：Goodreads、BookTok 与 Kindle 把暑期阅读从童年活动扩展为成年人的数字社交文化。
\end{itemize}

% ------------------------------------------------------------------------------
% 三、语调与韵律剖析
% ------------------------------------------------------------------------------
\readingsection{三、语调与韵律剖析}

\begin{prosodybox}[黄金句 1 · 意群与信息重音]
  \prosodyorig{There's nothing like a good book to help me escape when the weather is too hot and humid to enjoy doing anything else.}
  \prosodyipa{\textbf{humid} \ipa{/ˈhjuːmɪd/} \quad$\bm{\cdot}$\quad \textbf{escape} \ipa{/ɪˈskeɪp/}}
  \prosodychunks{There's \stress{NOTHING LIKE} a \stress{GOOD BOOK} \chunk to help me \stress{ESCAPE} \chunk when the weather is \stress{TOO HOT} and \stress{HUMID} \chunk to enjoy doing \stress{ANYTHING ELSE}. \tonefall}
  \prosodyfocus{nothing like, good book, escape, too hot, humid, anything else}
  \prosodyflaws{
    \item \texttt{humid} \ipa{/ˈhjuːmɪd/} 主重音在第一音节，避免误读为第二音节重音。
    \item \texttt{nothing like a} 不要逐词断开，整体保持平滑衔接。
  }
\end{prosodybox}

\begin{prosodybox}[黄金句 2 · 动作链与连贯输出]
  \prosodyorig{I eagerly awaited the release of the seventh Harry Potter book and bailed on family plans for two days in order to read it.}
  \prosodyipa{\textbf{eagerly} \ipa{/ˈiːɡəli/} \quad$\bm{\cdot}$\quad \textbf{bail} \ipa{/beɪl/}}
  \prosodychunks{I \stress{EAGERLY AWAITED} \chunk the release of the seventh \stress{HARRY POTTER} book \chunk and \stress{BAILED ON} family plans \chunk for \stress{TWO DAYS} \chunk \weak{in order to} \stress{READ IT}. \tonefall}
  \prosodyfocus{eagerly awaited, bailed on, two days, read it}
  \prosodyflaws{
    \item \phonlink{\texttt{bailed}}{\texttt{on}} 作为一个动作语块连贯说出，不在两个词之间人为停顿。
    \item \weak{in order to} 弱化并与后接动词迅速连带。
  }
\end{prosodybox}

% ------------------------------------------------------------------------------
% 四、输出训练
% ------------------------------------------------------------------------------
\readingsection{四、脱稿输出与观点}

\oralblock{快速概括}{
  In the US, summer reading began as an organized effort to keep children reading during a long school break. Over time, that habit has expanded into a broader digital reading culture through platforms such as Goodreads, BookTok and e-readers.
}

\oralblock{结构复述提纲}{
  \begin{itemize}
    \item 先解释为什么美国特别重视暑期阅读：暑假长，教育者担心 summer slide；
    \item 再讲图书馆怎样通过竞赛和奖励推动阅读；
    \item 接着用作者的童年经历做具体例子；
    \item 最后讲 Goodreads、BookTok、Kindle 怎样把这种文化延续到成年人的数字生活。
  \end{itemize}
}

\oralblock{个人观点}{
  数字平台未必只会削弱深度阅读。本文提供了一个相反案例：当社交机制、推荐系统和便携设备围绕阅读组织起来时，技术也可能让阅读行为更容易持续。
}

% ------------------------------------------------------------------------------
% 五、真实对谈：只在真实练习后记录
% ------------------------------------------------------------------------------
\readingsection{五、对谈与关键反馈}

\begin{aiinteraction}{为什么一种儿童暑期活动能够延伸成年人的数字阅读文化？}
  \aiquestion{Why do you think a childhood summer-reading habit can survive into adulthood in a digital form?}
  \aiuseranswer{【待真实对谈后记录学习者当时的原始回答。不要由 AI 代填。】}
  \aicorrection{
    \begin{itemize}
      \item 待真实回答后，只保留最值得修改的语言问题与目标表达调用情况。
    \end{itemize}
  }
  \ainatural{【仅在真实回答后，按需要补充自然改写。具有长期价值的新表达再进入表达层。】}
\end{aiinteraction}

% ------------------------------------------------------------------------------
% 六、本篇向外沉淀
% ------------------------------------------------------------------------------
\readingsection{六、本篇沉淀与跨层引用}

\begin{itemize}
  \item \textbf{\color{readGreen}已进入表达层}：\texttt{EXP-001} bail on；\texttt{EXP-002} counter the summer slide；\texttt{EXP-003} curl up in a cozy nook；\texttt{EXP-004} live up to the hype；\texttt{EXP-005} keep the magic of ... alive；\texttt{EXP-006} There's nothing like... to...；\texttt{EXP-007} amid the boom of...。
  \item \textbf{\color{readMain}仍留本篇}：highly organized effort、log books、summer slide、must-have accessory 等当前主要服务本文理解的词语。
  \item \textbf{\color{readGray}能力层锚点}：本篇可作为后续口语复述与目标表达调用的证据来源；只有真实训练完成后才更新长期能力结论。
\end{itemize}

\end{document}
